\section{Closure of the Probe--Closure Field System}
\begin{quote}\small
\noindent\textbf{Status.} \emph{Legacy field-closure block inside the active main body.}\par
\noindent\textbf{Block function.} This section closes the preceding field package and identifies the residue, limiting regimes, and Einstein-route consequences that remain readable afterward.
\end{quote}
\label{sec:closure-field-system}
%% ============================================================

\begin{definition}[Projected closure residue]
The local closure residue of the probe--closure system is defined by
\[
  \mathfrak R_C
  :=
  (\mathbf 1-\hat P_{C_7})\mathcal M_{16}\hat P_{C_7}
  + \hat P_{C_7}\mathcal M_{16}(\mathbf 1-\hat P_{C_7}).
\]
It measures the off-sector leakage of the Millennium matrix across the
Meta-Septum splitting.
The field system is called \emph{closed} if and only if
\[
  \mathfrak R_C = 0.
\]
\end{definition}

\begin{definition}[Explicit closure functional]
The compensator functional is taken in minimal form as
\[
  \Lambda[\mathcal M_{16},\hat P_{C_7}]
  :=
  \rho\,[\mathcal M_{16},\hat P_{C_7}]^
  \dagger [\mathcal M_{16},\hat P_{C_7}]
  + \sigma\,\mathfrak R_C^
  \dagger \mathfrak R_C,
\]
with effective coefficients $\rho,\sigma\ge 0$.
Hence
\[
  \mathcal H_C
  =
  \rho\,[\mathcal M_{16},\hat P_{C_7}]^
  \dagger [\mathcal M_{16},\hat P_{C_7}]
  + \sigma\,\mathfrak R_C^
  \dagger \mathfrak R_C.
\]
This realizes the compensator as a positive semidefinite measure of
closure mismatch.
\end{definition}

\begin{proposition}[Equivalent closure criteria]
\label{prop:equivalent-closure-criteria}
The following statements are equivalent in the effective probe sector:
\begin{enumerate}
  \item $\mathfrak R_C=0$.
  \item $[\mathcal M_{16},\hat P_{C_7}]=0$.
  \item $\mathcal M_{16}$ preserves the $C_7$-sector and its complement.
\end{enumerate}
In particular, exact closure means that the Millennium matrix becomes
block-diagonal with respect to the decomposition
\[
  \mathcal H_{\mathrm{sept}}
  = \hat P_{C_7}\mathcal H_{\mathrm{sept}}
  \oplus
  (\mathbf 1-\hat P_{C_7})\mathcal H_{\mathrm{sept}}.
\]
\end{proposition}

\begin{proof}
Writing $Q:=\hat P_{C_7}$ and $Q_\perp:=\mathbf 1-Q$, we have
\[
  [\mathcal M_{16},Q]
  =
  Q_\perp\mathcal M_{16}Q - Q\mathcal M_{16}Q_\perp.
\]
Hence $[\mathcal M_{16},Q]=0$ if and only if both mixed blocks vanish,
which is exactly the condition $\mathfrak R_C=0$.
This is equivalent to invariance of the two summands under
$\mathcal M_{16}$, i.e. block-diagonality in the stated splitting.
\end{proof}

\begin{theorem}[Closed master system]
\label{thm:closed-master-system}
The probe--closure field equations are internally closed by the system
\begin{align*}
  \bigl(i\Gamma^\mu\mathcal D_\mu - \Omega\bigr)\Psi &= 0,\\
  \mathcal D^\mu\mathcal F_{\mu\nu}
    &= J^{(\Psi)}_\nu + J^{(C_7)}_\nu + J^{(HC)}_\nu,\\
  (\mathbf 1-\hat P_{C_7})\Psi &= 0,\\
  \mathcal H_C &= \Lambda[\mathcal M_{16},\hat P_{C_7}],\\
  \mathfrak R_C &= 0.
\end{align*}
Therefore all previously free closure symbols are eliminated in favour
of explicit algebraic conditions on $(\mathcal M_{16},\hat P_{C_7})$.
\end{theorem}

\begin{corollary}[Vanishing compensator in the exact sector]
If the closure condition $\mathfrak R_C=0$ holds exactly, then
\[
  [\mathcal M_{16},\hat P_{C_7}] = 0
  \qquad\text{and hence}\qquad
  \mathcal H_C = 0
\]
for the minimal compensator ansatz above.
Thus the Heisenberg compensator is nonzero only away from exact closure,
where it measures and suppresses leakage out of the stable sector.
\end{corollary}

\begin{remark}[Interpretation of the closure step]
The previous section produced a coupled field architecture.
The present section closes it in the minimal algebraic sense:
physical states stay in the $C_7$-sector, the Millennium matrix may not
mix stable and unstable blocks, and the compensator becomes an explicit
positive residue functional rather than a symbolic placeholder.
This does not yet derive gravitation or the Standard Model, but it does
remove the main formal gap inside the probe sector.
\end{remark}

\begin{quote}\small
\noindent\textbf{Reader cue.} Read the next three stages as a layered unpacking of the same closure statement. Its purpose is to move from the core algebra to leakage/confinement language, then to controlled limiting regimes, and finally to the explicit Einstein-route reconstruction.
\end{quote}

\subsection*{Physical interpretation of the closure residue}

\begin{remark}[Leakage, confinement, and stability]
The residue $\mathfrak R_C$ has a direct physical reading.
Its mixed blocks measure transitions between the stable $C_7$-sector and
its complement.
Hence $\mathfrak R_C\neq 0$ represents \,leakage\, of probe amplitude out
of the metastably confined sector, while $\mathfrak R_C=0$ expresses exact
sector confinement.
In this sense the closure condition is the algebraic analogue of a
superselection rule: admissible long-lived states are precisely those for
which the effective internal dynamics generated by $\mathcal M_{16}$ never
pushes the state out of the projected septum sector.
\end{remark}

\begin{remark}[Role of the compensator]
For the minimal ansatz,
\[
  \mathcal H_C
  = \rho\,[\mathcal M_{16},\hat P_{C_7}]^\dagger[\mathcal M_{16},\hat P_{C_7}]
    + \sigma\,\mathfrak R_C^\dagger\mathfrak R_C,
\]
the compensator is not an independent interaction but an energetic penalty
for sector mixing.
Its effect is twofold:
first, it raises the cost of off-sector propagation; second, it drives the
variation toward block preservation of $\mathcal M_{16}$.
Thus the compensator plays the role of a restoring term that suppresses
unstable excursions and selects dynamically persistent probe states.
\end{remark}

\begin{remark}[Physical reading of exact closure]
In the exactly closed regime the splitting
\[
  \mathcal H_{\mathrm{sept}}
  = \hat P_{C_7}\mathcal H_{\mathrm{sept}}
  \oplus
  (\mathbf 1-\hat P_{C_7})\mathcal H_{\mathrm{sept}}
\]
is preserved by the internal generator $\mathcal M_{16}$.
Then the effective dynamics on the physical sector becomes autonomous.
Operationally this means that all observables built from $\Psi$ may be
computed within the projected sector without referring to hidden leakage
channels.
The closed system is therefore the minimal regime in which the probe theory
admits a stable phenomenological interpretation.
\end{remark}

\subsection*{Controlled limiting regimes}

\begin{proposition}[Dirac limit]
\label{prop:dirac-limit}
Assume the exact closure condition $\mathfrak R_C=0$, freeze the connection
background to a locally trivial configuration $\mathcal A_\mu=0$, and take
$\Omega$ to reduce on the closed sector to a scalar effective mass,
\[
  \Omega\big|_{\hat P_{C_7}\mathcal H}
  = m_{\mathrm{eff}}\,\mathbf 1.
\]
Then the master equation reduces on the physical sector to
\[
  (i\Gamma^\mu\partial_\mu - m_{\mathrm{eff}})\Psi = 0,
\]
which is the ordinary free Dirac form up to the choice of the effective
gamma representation.
\end{proposition}

\begin{proof}
Under exact closure the constraint $(\mathbf 1-\hat P_{C_7})\Psi=0$ removes
off-sector components, and for the minimal ansatz the compensator vanishes.
With $\mathcal A_\mu=0$ and scalar reduction of $\Omega$, the covariant
operator collapses to $\partial_\mu$, so the master system becomes the free
Dirac equation on the projected sector.
\end{proof}

\begin{proposition}[Yang--Mills type limit]
\label{prop:yang-mills-limit}
Assume again exact closure, but retain a nontrivial connection
$\mathcal A_\mu$ acting internally on the closed sector.
If the projected operator algebra generated by $\mathcal A_\mu$ closes on a
finite-dimensional Lie algebra $\mathfrak g$, then the resonance field
equation
\[
  \mathcal D^\mu \mathcal F_{\mu\nu}=J_\nu
\]
reduces on the physical sector to a Yang--Mills type equation with gauge
algebra $\mathfrak g$.
\end{proposition}

\begin{proof}
Once $\hat P_{C_7}$ is preserved by $\mathcal M_{16}$, the connection and its
curvature may be restricted consistently to the closed sector.
If the projected commutator algebra is Lie-closed, then the curvature takes
the standard form of a gauge field strength on $\mathfrak g$, and the field
equation is exactly of Yang--Mills type with source $J_\nu$.
\end{proof}

\begin{proposition}[Einstein-type effective limit]
\label{prop:einstein-effective-limit}
Suppose the projected connection separates into an external spacetime part
and an internal resonance part,
\[
  \mathcal D_\mu = \nabla_\mu^{(g)} + \mathbb A_\mu,
\]
where $\nabla_\mu^{(g)}$ is compatible with a coarse-grained metric
$g_{\mu\nu}$ extracted from bilinear probe observables.
If the effective action is integrated over fast internal modes and the
leading derivative expansion is local, then the bosonic sector has the
generic low-energy form
\[
  S_{\mathrm{eff}}[g,\mathbb A]
  = \int d^4x\sqrt{|g|}\,
    \bigl(c_0 + c_1 R(g) + c_2\operatorname{Tr}(\mathbb F_{\mu\nu}\mathbb F^{\mu\nu}) + \cdots\bigr).
\]
Hence the metric sector obeys Einstein-type equations at leading order.
\end{proposition}

\begin{remark}[Status of the Einstein limit]
This is presently a controlled effective-theory statement, not yet a first-
principles derivation.
What closure provides is the prerequisite for such a limit: a dynamically
stable projected sector, a well-defined connection, and a variationally
closed action.
To obtain actual Einstein equations one still has to specify how the metric
is reconstructed from probe bilinears and how the internal modes are coarse-
grained.
\end{remark}

\begin{quote}\small
\noindent\textbf{Reader cue.} Read the next block as a concrete reconstruction route for the Einstein-type limit. Its purpose is to spell out the bilinear coframe, the induced metric, and the coarse-grained assumptions under which the Einstein tensor reappears.
\end{quote}

\subsection*{A concrete route to the Einstein limit}

We now make the effective Einstein-type limit more explicit at the level of
field reconstruction. Let $\Psi_\triangle=(\Psi_1,\Psi_2,\Psi_\perp)^T$ be the
closed triadic probe field and let the exact closure condition
$\mathfrak R_C=0$ hold. On the physical sector we define the bilinear one-forms
\[
  E^a{}_{\mu}
  := \operatorname{Re}\bigl(\bar\Psi_\triangle \,\Gamma^a \, \mathcal D_\mu
  \Psi_\triangle\bigr),
  \qquad a=0,1,2,3,
\]
where $\Gamma^a$ denotes a fixed Clifford frame acting on the projected probe
space. Whenever the $4\times 4$ matrix $E^a{}_{\mu}$ has maximal rank, these
bilinears define an effective coframe and hence an effective metric
\[
  g_{\mu\nu}
  := \eta_{ab} E^a{}_{\mu} E^b{}_{\nu},
\]
with $\eta_{ab}=\operatorname{diag}(-1,1,1,1)$.

The closure constraint is essential here: it ensures that the bilinears are
computed entirely within the invariant $C_7$-sector, so the reconstructed
metric is not contaminated by leakage terms. In this sense $g_{\mu\nu}$ is not
introduced externally but induced from the closed probe dynamics.

Next, restrict the covariant operator to the physical sector,
\[
  \mathcal D^{(C)}_\mu := \hat P_{C_7}\mathcal D_\mu \hat P_{C_7}.
\]
Assume that $\mathcal D^{(C)}_\mu$ admits a decomposition into a geometric and an
internal part,
\[
  \mathcal D^{(C)}_\mu
  = \partial_\mu + \frac14 \, \omega_{\mu ab}\,\Sigma^{ab} + \mathbb A_\mu,
\]
where $\Sigma^{ab}=\frac12[\Gamma^a,\Gamma^b]$ and $\omega_{\mu ab}$ is chosen so
that the induced connection is metric-compatible,
\[
  \nabla^{(g)}_\lambda g_{\mu\nu}=0.
\]
If in addition the torsion-free condition is imposed at the coarse-grained
level, then $\omega_{\mu ab}$ coincides with the Levi--Civita spin connection of
$g_{\mu\nu}$.

This yields the following sharpened version of the Einstein limit.

\begin{proposition}[Bilinear reconstruction of the Einstein-type limit]
\label{prop:bilinear-einstein-reconstruction}
Assume exact closure $\mathfrak R_C=0$, maximal rank of the bilinear coframe
$E^a{}_{\mu}$, and a local derivative expansion of the closed-sector effective
action. Suppose moreover that the coarse-grained bosonic action takes the form
\[
  S_{\mathrm{eff}}[g,\mathbb A,\Phi]
  = \int d^4x\sqrt{|g|}\,\bigl(
      \Lambda_{\mathrm{eff}}
      + \frac{1}{2\kappa_{\mathrm{eff}}}R(g)
      - \frac14\operatorname{Tr}(\mathbb F_{\mu\nu}\mathbb F^{\mu\nu})
      + \mathcal L_{\mathrm{cl}}(\Phi,g)
    \bigr),
\]
where $\Phi$ denotes the remaining closure-sector collective fields.
Then variation with respect to $g^{\mu\nu}$ yields the effective equation
\[
  G_{\mu\nu}(g) + \Lambda_{\mathrm{eff}}g_{\mu\nu}
  = \kappa_{\mathrm{eff}}
    \bigl(T^{(\mathbb A)}_{\mu\nu} + T^{(\Phi)}_{\mu\nu}\bigr),
\]
with the standard stress-energy tensors obtained from the gauge and collective
sectors. In particular, the Einstein tensor arises from the closed probe
system once the metric is reconstructed from bilinears and the internal modes
are integrated out.
\end{proposition}

\begin{proof}
Under exact closure the projected operator $\mathcal D^{(C)}_\mu$ acts
autonomously on the physical sector. Maximal rank of $E^a{}_{\mu}$ turns the
probe bilinears into a nondegenerate coframe, hence into a metric
$g_{\mu\nu}$. Metric compatibility and vanishing torsion identify the geometric
part of $\mathcal D^{(C)}_\mu$ with the Levi--Civita spin connection.
By assumption, coarse-graining over the fast internal modes produces a local
effective action whose leading geometric term is the Einstein--Hilbert term
$R(g)$. Varying $S_{\mathrm{eff}}$ with respect to $g^{\mu\nu}$ gives the Einstein
tensor on the left-hand side and the stress-energy tensors of the remaining
fields on the right-hand side, which is the claimed Einstein-type equation.
\end{proof}

\begin{remark}[What is still model-dependent]
The above construction isolates the precise points where additional work is
still required.
First, one must justify the nondegeneracy of the bilinear coframe from the
triadic probe dynamics rather than postulate it.
Second, one has to derive the constants $\Lambda_{\mathrm{eff}}$ and
$\kappa_{\mathrm{eff}}$ from the microscopic closure model.
Third, the collective closure fields $\Phi$ have to be identified explicitly;
they encode the backreaction of the residual probe, projector, and compensator
degrees of freedom on spacetime geometry.
What is gained already is structural: the Einstein sector now appears by a
clear reconstruction pipeline
\[
  (\Psi_\triangle,\hat P_{C_7},\mathcal H_C,\mathcal M_{16})
  \longrightarrow E^a{}_{\mu}
  \longrightarrow g_{\mu\nu}
  \longrightarrow S_{\mathrm{eff}}[g,\mathbb A,\Phi]
  \longrightarrow G_{\mu\nu}=\kappa_{\mathrm{eff}}T_{\mu\nu}^{\mathrm{eff}}.
\]
\end{remark}

\begin{remark}[Interpretive summary]
The probe--closure system now supports a three-level reading.
At the exact algebraic level, closure means vanishing leakage between the
physical $C_7$-sector and its complement.
At the dynamical level, the compensator suppresses deviations from that
sector.
At the phenomenological level, ordinary Dirac and Yang--Mills equations
appear as closed-sector limits, while Einstein dynamics can arise as a
further coarse-grained low-energy description if a metric reconstruction is
available.
\end{remark}

\begin{quote}\small
\noindent\textbf{Historical transition note.} Read the next synthesis block as archival proof-memory. Its purpose is to preserve back-reference, genealogy, and an earlier integrative summary beneath the active main line, not to add a second live closure engine.
\end{quote}

%% ============================================================



